\documentclass[10pt,a4paper]{amsart}
\usepackage{amsmath,amssymb,amsthm,graphicx}
% Retained minimal preamble for the Stacks Project chapter sites-cohomology.tex (Cohomology on Sites).
% Verbatim excerpt of the chapter preamble: Xy-pic setup (native preamble lines 14-18) and the native
% theorem declarations with the mathematical macros (native preamble lines 154-217). Copied mechanically;
% nothing added, removed or redefined. The wrapper class and the amsmath/amssymb/amsthm/graphicx base
% packages are supplied by the assembler, so the class/externaldocument part of the native preamble is omitted.
% For commutative diagrams we use Xy-pic
\usepackage[all]{xy}

% We use 2cell for 2-commutative diagrams.
\xyoption{2cell}
% Theorem environments.
%
\theoremstyle{plain}
\newtheorem{theorem}[subsection]{Theorem}
\newtheorem{proposition}[subsection]{Proposition}
\newtheorem{lemma}[subsection]{Lemma}

\theoremstyle{definition}
\newtheorem{definition}[subsection]{Definition}
\newtheorem{example}[subsection]{Example}
\newtheorem{exercise}[subsection]{Exercise}
\newtheorem{situation}[subsection]{Situation}

\theoremstyle{remark}
\newtheorem{remark}[subsection]{Remark}
\newtheorem{remarks}[subsection]{Remarks}

\numberwithin{equation}{subsection}

% Macros
%
\def\lim{\mathop{\mathrm{lim}}\nolimits}
\def\colim{\mathop{\mathrm{colim}}\nolimits}
\def\Spec{\mathop{\mathrm{Spec}}}
\def\Hom{\mathop{\mathrm{Hom}}\nolimits}
\def\Ext{\mathop{\mathrm{Ext}}\nolimits}
\def\SheafHom{\mathop{\mathcal{H}\!\mathit{om}}\nolimits}
\def\SheafExt{\mathop{\mathcal{E}\!\mathit{xt}}\nolimits}
\def\Sch{\mathit{Sch}}
\def\Mor{\mathop{\mathrm{Mor}}\nolimits}
\def\Ob{\mathop{\mathrm{Ob}}\nolimits}
\def\Sh{\mathop{\mathit{Sh}}\nolimits}
\def\NL{\mathop{N\!L}\nolimits}
\def\CH{\mathop{\mathrm{CH}}\nolimits}
\def\proetale{{pro\text{-}\acute{e}tale}}
\def\etale{{\acute{e}tale}}
\def\QCoh{\mathit{QCoh}}
\def\Ker{\mathop{\mathrm{Ker}}}
\def\Im{\mathop{\mathrm{Im}}}
\def\Coker{\mathop{\mathrm{Coker}}}
\def\Coim{\mathop{\mathrm{Coim}}}

% Boxtimes
%
\DeclareMathSymbol{\boxtimes}{\mathbin}{AMSa}{"02}

%
% Macros for moduli stacks/spaces
%
\def\QCohstack{\mathcal{QC}\!\mathit{oh}}
\def\Cohstack{\mathcal{C}\!\mathit{oh}}
\def\Spacesstack{\mathcal{S}\!\mathit{paces}}
\def\Quotfunctor{\mathrm{Quot}}
\def\Hilbfunctor{\mathrm{Hilb}}
\def\Curvesstack{\mathcal{C}\!\mathit{urves}}
\def\Polarizedstack{\mathcal{P}\!\mathit{olarized}}
\def\Complexesstack{\mathcal{C}\!\mathit{omplexes}}
% \Pic is the operator that assigns to X its picard group, usage \Pic(X)
% \Picardstack_{X/B} denotes the Picard stack of X over B
% \Picardfunctor_{X/B} denotes the Picard functor of X over B
\def\Pic{\mathop{\mathrm{Pic}}\nolimits}
\def\Picardstack{\mathcal{P}\!\mathit{ic}}
\def\Picardfunctor{\mathrm{Pic}}
\def\Deformationcategory{\mathcal{D}\!\mathit{ef}}
\usepackage{cleveref}
\crefname{originalsection}{original section}{original sections}
\Crefname{originalsection}{Original section}{Original sections}
\makeatletter
\newlabel{section-compare}{{0EWK}{source}{}{sourcepointer-1}{}}
\newlabel{section-compare@cref}{{[originalsection][0][]0EWK}{[1][0][]source}{}{}{}}
\makeatother
\crefname{originallemma}{original lemma}{original lemmas}
\Crefname{originallemma}{Original lemma}{Original lemmas}
\makeatletter
\newlabel{lemma-higher-direct-images}{{072W}{source}{}{sourcepointer-2}{}}
\newlabel{lemma-higher-direct-images@cref}{{[originallemma][0][]072W}{[1][0][]source}{}{}{}}
\makeatother
\crefname{originallemma}{original lemma}{original lemmas}
\Crefname{originallemma}{Original lemma}{Original lemmas}
\makeatletter
\newlabel{lemma-kill-cohomology-class-on-covering}{{01FW}{source}{}{sourcepointer-3}{}}
\newlabel{lemma-kill-cohomology-class-on-covering@cref}{{[originallemma][0][]01FW}{[1][0][]source}{}{}{}}
\makeatother
\crefname{originallemma}{original lemma}{original lemmas}
\Crefname{originallemma}{Original lemma}{Original lemmas}
\makeatletter
\newlabel{lemma-V-implies-cohomology-extra-general}{{0EZC}{source}{}{sourcepointer-4}{}}
\newlabel{lemma-V-implies-cohomology-extra-general@cref}{{[originallemma][0][]0EZC}{[1][0][]source}{}{}{}}
\makeatother
\crefname{originallemma}{original lemma}{original lemmas}
\Crefname{originallemma}{Original lemma}{Original lemmas}
\makeatletter
\newlabel{lemma-make-class-zero-general}{{0EZD}{source}{}{sourcepointer-5}{}}
\newlabel{lemma-make-class-zero-general@cref}{{[originallemma][0][]0EZD}{[1][0][]source}{}{}{}}
\makeatother
\crefname{originallemma}{original lemma}{original lemmas}
\Crefname{originallemma}{Original lemma}{Original lemmas}
\makeatletter
\newlabel{lemma-Leray-unbounded}{{0D6H}{source}{}{sourcepointer-6}{}}
\newlabel{lemma-Leray-unbounded@cref}{{[originallemma][0][]0D6H}{[1][0][]source}{}{}{}}
\makeatother
\crefname{originallemma}{original lemma}{original lemmas}
\Crefname{originallemma}{Original lemma}{Original lemmas}
\makeatletter
\newlabel{derived-remark-truncation-distinguished-triangle}{{08J5}{source}{}{sourcepointer-7}{}}
\newlabel{derived-remark-truncation-distinguished-triangle@cref}{{[originallemma][0][]08J5}{[1][0][]source}{}{}{}}
\makeatother
\crefname{originalequation}{original equation}{original equations}
\Crefname{originalequation}{Original equation}{Original equations}
\makeatletter
\newlabel{equation-commutative-epsilon}{{0EZ2}{source}{}{sourcepointer-8}{}}
\newlabel{equation-commutative-epsilon@cref}{{[originalequation][0][]0EZ2}{[1][0][]source}{}{}{}}
\makeatother
\crefname{originallemma}{original lemma}{original lemmas}
\Crefname{originallemma}{Original lemma}{Original lemmas}
\makeatletter
\newlabel{lemma-V-implies-C-general}{{0EZA}{source}{}{sourcepointer-9}{}}
\newlabel{lemma-V-implies-C-general@cref}{{[originallemma][0][]0EZA}{[1][0][]source}{}{}{}}
\makeatother
\crefname{originallemma}{original lemma}{original lemmas}
\Crefname{originallemma}{Original lemma}{Original lemmas}
\makeatletter
\newlabel{lemma-A}{{0EZ9}{source}{}{sourcepointer-10}{}}
\newlabel{lemma-A@cref}{{[originallemma][0][]0EZ9}{[1][0][]source}{}{}{}}
\makeatother
\crefname{originallemma}{original lemma}{original lemmas}
\Crefname{originallemma}{Original lemma}{Original lemmas}
\makeatletter
\newlabel{lemma-V-implies-cohomology-general}{{0EZB}{source}{}{sourcepointer-11}{}}
\newlabel{lemma-V-implies-cohomology-general@cref}{{[originallemma][0][]0EZB}{[1][0][]source}{}{}{}}
\makeatother
\title{The induction step from property V-n to property V-n-plus-one in the comparison of topologies}
\author{The Stacks Project Authors}
\date{stacks-project:sites-cohomology.tex:local-snapshot-20261011; GNU FDL 1.2 or later}
\begin{document}
\maketitle
\noindent\textit{Editorial note.} All delivered mathematical text is the author's own, copied from the pinned Stacks Project chapter file sites-cohomology.tex; no statement or proof step was rewritten, reordered or removed.\par\medskip
\noindent\textit{Editorial note.} The counted statement reads 'In Situation ... we have (V\_n) implies (V\_\{n+1\})', so the situation environment of the comparison section is delivered as a context span (source lines 6840-6879); it is the author's own text copied verbatim and makes the card self-contained.\par\medskip
\noindent\textit{Editorial note.} The proof is the author's complete argument: the four stages above are all present inside the counted span, and the closing reduction is carried through rather than summarised.\par\medskip
\noindent\textit{Editorial note.} Every labelled reference occurring inside the counted spans points to a chapter of the Project that is not part of this card's source tree. Each is carried in assembly.external\_references as a non-load-bearing navigation pointer whose source location is the invocation inside the counted span, so no load-bearing mathematics is replaced by a stub.\par\medskip
\noindent\textit{Editorial note.} The preamble is the chapter's own retained minimal preamble (sources/stacks-sites-cohomology/retained-preamble.tex), a verbatim excerpt of the native preamble, which declares the situation environment used by the context span.\par\medskip
\noindent\textit{Allowed standard prerequisites (editorial).}\begin{itemize}
\item The situation of the comparison section: the site, the two topologies tau and tau', the set of arrows P and the weak Serre subcategories A'\_X. Same chapter, the situation environment of the comparison section, delivered verbatim as a context span in this card (source lines 6840-6879). Supplies every object and hypothesis occurring in the statement and in the proof.
\item Property (V\_n) of the comparison. Same chapter, the property (V\_n) stating that every tau'-cohomology class of degree n on an object of C dies after restricting to the members of a tau'-covering and pulling back along an arrow of P. Supplies the hypothesis from which (V\_\{n+1\}) is deduced.
\item Locality of cohomology and the cohomology comparison in degree n. Same chapter, the locality lemma for cohomology classes, the comparison result identifying tau'-classes on X after pullback along an arrow of P, and the Leray-type lemma used at the end. Supplies the two reductions at the beginning of the proof and the final computation with the distinguished triangle attached to F -> Rf\_\{tau,*\}f\_tau\textasciicircum{}\{-1\}F.
\end{itemize}\medskip
\ifdefined\hypertarget\hypertarget{sourcepointer-1}{}\fi
\noindent\textit{Source reference.} Original section 0EWK: sources/stacks-sites-cohomology/pinned-original-source.tex, source lines 6177--6177. Navigation pointer invoked at the counted span; the counted span contains the pointer only and the cited result is not reproduced inside it, so the pointer is registered as a non-load-bearing navigation stub in the established form for Stacks material. No bibliography citation is involved.\par\medskip
\ifdefined\hypertarget\hypertarget{sourcepointer-2}{}\fi
\noindent\textit{Source reference.} Original lemma 072W: sources/stacks-sites-cohomology/pinned-original-source.tex, source lines 747--747. Navigation pointer invoked at proof line 7157; the counted span contains the pointer only and the cited result is not reproduced inside it, so the pointer is registered as a non-load-bearing navigation stub in the established form for Stacks material. No bibliography citation is involved.\par\medskip
\ifdefined\hypertarget\hypertarget{sourcepointer-3}{}\fi
\noindent\textit{Source reference.} Original lemma 01FW: sources/stacks-sites-cohomology/pinned-original-source.tex, source lines 715--715. Navigation pointer invoked at proof line 7163; the counted span contains the pointer only and the cited result is not reproduced inside it, so the pointer is registered as a non-load-bearing navigation stub in the established form for Stacks material. No bibliography citation is involved.\par\medskip
\ifdefined\hypertarget\hypertarget{sourcepointer-4}{}\fi
\noindent\textit{Source reference.} Original lemma 0EZC: sources/stacks-sites-cohomology/pinned-original-source.tex, source lines 7061--7061. Navigation pointer invoked at proof line 7184; the counted span contains the pointer only and the cited result is not reproduced inside it, so the pointer is registered as a non-load-bearing navigation stub in the established form for Stacks material. No bibliography citation is involved.\par\medskip
\ifdefined\hypertarget\hypertarget{sourcepointer-5}{}\fi
\noindent\textit{Source reference.} Original lemma 0EZD: sources/stacks-sites-cohomology/pinned-original-source.tex, source lines 7090--7090. Navigation pointer invoked at proof line 7191; the counted span contains the pointer only and the cited result is not reproduced inside it, so the pointer is registered as a non-load-bearing navigation stub in the established form for Stacks material. No bibliography citation is involved.\par\medskip
\ifdefined\hypertarget\hypertarget{sourcepointer-6}{}\fi
\noindent\textit{Source reference.} Original lemma 0D6H: sources/stacks-sites-cohomology/pinned-original-source.tex, source lines 4403--4403. Navigation pointer invoked at proof line 7209; the counted span contains the pointer only and the cited result is not reproduced inside it, so the pointer is registered as a non-load-bearing navigation stub in the established form for Stacks material. No bibliography citation is involved.\par\medskip
\ifdefined\hypertarget\hypertarget{sourcepointer-7}{}\fi
\noindent\textit{Source reference.} Original lemma 08J5: sources/stacks-sites-cohomology/pinned-original-source.tex, source lines 7212--7212. Navigation pointer invoked at proof line 7212; the counted span contains the pointer only and the cited result is not reproduced inside it, so the pointer is registered as a non-load-bearing navigation stub in the established form for Stacks material. No bibliography citation is involved.\par\medskip
\ifdefined\hypertarget\hypertarget{sourcepointer-8}{}\fi
\noindent\textit{Source reference.} Original equation 0EZ2: sources/stacks-sites-cohomology/pinned-original-source.tex, source lines 6817--6817. Navigation pointer invoked at proof line 7225; the counted span contains the pointer only and the cited result is not reproduced inside it, so the pointer is registered as a non-load-bearing navigation stub in the established form for Stacks material. No bibliography citation is involved.\par\medskip
\ifdefined\hypertarget\hypertarget{sourcepointer-9}{}\fi
\noindent\textit{Source reference.} Original lemma 0EZA: sources/stacks-sites-cohomology/pinned-original-source.tex, source lines 6914--6914. Navigation pointer invoked at proof line 7231; the counted span contains the pointer only and the cited result is not reproduced inside it, so the pointer is registered as a non-load-bearing navigation stub in the established form for Stacks material. No bibliography citation is involved.\par\medskip
\ifdefined\hypertarget\hypertarget{sourcepointer-10}{}\fi
\noindent\textit{Source reference.} Original lemma 0EZ9: sources/stacks-sites-cohomology/pinned-original-source.tex, source lines 6882--6882. Navigation pointer invoked at proof line 7245; the counted span contains the pointer only and the cited result is not reproduced inside it, so the pointer is registered as a non-load-bearing navigation stub in the established form for Stacks material. No bibliography citation is involved.\par\medskip
\ifdefined\hypertarget\hypertarget{sourcepointer-11}{}\fi
\noindent\textit{Source reference.} Original lemma 0EZB: sources/stacks-sites-cohomology/pinned-original-source.tex, source lines 7029--7029. Navigation pointer invoked at proof line 7275; the counted span contains the pointer only and the cited result is not reproduced inside it, so the pointer is registered as a non-load-bearing navigation stub in the established form for Stacks material. No bibliography citation is involved.\par\medskip
\begin{situation}
\label{situation-compare}
With $\mathcal{C}$, $\tau$, and $\tau'$ as in Section \ref{section-compare}.
Assume we are given a subset $\mathcal{P} \subset \text{Arrows}(\mathcal{C})$
and for every object $X$ of $\mathcal{C}$ we are given a weak Serre subcategory
$\mathcal{A}'_X \subset \textit{Ab}(\mathcal{C}_{\tau'}/X)$.
We make the following assumption:
\begin{enumerate}
\item
\label{item-base-change-P}
given $f : X \to Y$ in $\mathcal{P}$ and $Y' \to Y$ general,
then $X \times_Y Y'$ exists and $X \times_Y Y' \to Y'$ is in $\mathcal{P}$,
\item
\label{item-restriction-A}
$f_{\tau'}^{-1}$ sends $\mathcal{A}'_Y$ into $\mathcal{A}'_X$
for any morphism $f : X \to Y$ of $\mathcal{C}$,
\item
\label{item-A-sheaf}
given $X$ in $\mathcal{C}$ and $\mathcal{F}'$ in $\mathcal{A}'_X$, then
$\mathcal{F}'$ satisfies the sheaf condition for $\tau$-coverings, i.e.,
$\mathcal{F}' = \epsilon_{X, *}\epsilon_X^{-1}\mathcal{F}'$,
\item
\label{item-A-and-P}
if $f : X \to Y$ in $\mathcal{P}$ and
$\mathcal{F}' \in \Ob(\mathcal{A}'_X)$, then
$R^if_{\tau', *}\mathcal{F}' \in \Ob(\mathcal{A}'_Y)$
for $i \geq 0$.
\item
\label{item-refine-tau-by-P}
if $\{U_i \to U\}_{i \in I}$ is a $\tau$-covering, then there exist
\begin{enumerate}
\item a $\tau'$-covering $\{V_j \to U\}_{j \in J}$,
\item a $\tau$-covering $\{f_j : W_j \to V_j\}$ consisting
of a single $f_j \in \mathcal{P}$, and
\item a $\tau'$-covering $\{W_{jk} \to W_j\}_{k \in K_j}$
\end{enumerate}
such that $\{W_{jk} \to U\}_{j \in J, k \in K_j}$ is a refinement of
$\{U_i \to U\}_{i \in I}$.
\end{enumerate}
\end{situation}

\begin{lemma}
\label{lemma-induction-step-V-C-general}
In Situation \ref{situation-compare} we have $(V_n) \Rightarrow (V_{n + 1})$.
\end{lemma}

\begin{proof}
Let $X$ in $\mathcal{C}$ and $\mathcal{F}$ in $\mathcal{A}_X$.
Let $\xi \in H^{n + 1}_\tau(U, \mathcal{F})$ for some $U/X$.
We have to show that $\xi$ restricts to zero on the members of
a $\tau'$-covering of $U$. See
Lemma \ref{lemma-higher-direct-images}.
It follows from this that we may
replace $U$ by the members of a $\tau'$-covering of $U$.

\medskip\noindent
By locality of cohomology
(Lemma \ref{lemma-kill-cohomology-class-on-covering})
we can choose a $\tau$-covering $\{U_i \to U\}$
such that $\xi$ restricts to zero on $U_i$.
Choose $\{V_j \to V\}$, $\{f_j : W_j \to V_j\}$, and
$\{W_{jk} \to W_j\}$ as in (\ref{item-refine-tau-by-P}).
After replacing both $U$ by $V_j$
and $\mathcal{F}$ by its restriction to
$\mathcal{C}_\tau/V_j$, which is allowed by
(\ref{item-base-change-P}), we reduce to the
case discussed in the next paragraph.

\medskip\noindent
Here $f : X \to Y$ is an element of $\mathcal{P}$
such that $\{X \to Y\}$ is a $\tau$-covering,
$\mathcal{F}$ is an object of $\mathcal{A}_Y$, and
$\xi \in H^{n + 1}_\tau(Y, \mathcal{F})$ is such that
there exists a $\tau'$-covering $\{X_i \to X\}_{i \in I}$
such that $\xi$ restricts to zero on $X_i$ for all $i \in I$.
Problem: show that $\xi$ restricts to zero on a $\tau'$-covering of $Y$.

\medskip\noindent
By Lemma \ref{lemma-V-implies-cohomology-extra-general} there exists a
unique $\tau'$-cohomology class
$\theta \in H^{n + 1}_{\tau'}(X, \epsilon_{X, *}\mathcal{F})$
whose image is $\xi|_X$.
Since $\xi|_X$ pulls back to the same class on $X \times_Y X$
via the two projections, we find that the same is true for $\theta$
(by uniqueness).
By Lemma \ref{lemma-make-class-zero-general}
we see that after replacing $Y$ by the members of a $\tau'$-covering,
we may assume that $\theta = 0$.
Consequently, we may assume that $\xi|_X$ is zero.

\medskip\noindent
Let $f : X \to Y$ be an element of $\mathcal{P}$
such that $\{X \to Y\}$ is a $\tau$-covering,
$\mathcal{F}$ is an object of $\mathcal{A}_Y$, and
$\xi \in H^{n + 1}_\tau(Y, \mathcal{F})$
maps to zero in $H^{n + 1}_\tau(X, \mathcal{F})$.
Problem: show that $\xi$ restricts to zero on a $\tau'$-covering of $Y$.

\medskip\noindent
The assumptions tell us $\xi$ maps to zero under the map
$$
\mathcal{F} \longrightarrow Rf_{\tau, *}f_\tau^{-1}\mathcal{F}
$$
Use Lemma \ref{lemma-Leray-unbounded}.
A simple argument using the distinguished triangle of truncations
(Derived Categories, Remark
\ref{derived-remark-truncation-distinguished-triangle}) shows that
$\xi$ maps to zero under the map
$$
\mathcal{F} \longrightarrow \tau_{\leq n}Rf_{\tau, *}f_\tau^{-1}\mathcal{F}
$$
We will compare this with the map $\epsilon_{Y, *}\mathcal{F} \to K$
where
$$
K = \tau_{\leq n}Rf_{\tau', *}f_{\tau'}^{-1}\epsilon_{Y, *}\mathcal{F} =
\tau_{\leq n}Rf_{\tau', *}\epsilon_{X, *}f_{\tau}^{-1}\mathcal{F}
$$
The equality
$\epsilon_{X, *} f_\tau^{-1} = f_{\tau'}^{-1} \epsilon_{Y, *}$
is a property of (\ref{equation-commutative-epsilon}). Consider the map
$$
Rf_{\tau', *}\epsilon_{X, *}f_{\tau}^{-1}\mathcal{F} \longrightarrow
Rf_{\tau', *}R\epsilon_{X, *}f_{\tau}^{-1}\mathcal{F} =
R\epsilon_{Y, *}Rf_{\tau, *}f_\tau^{-1}\mathcal{F}
$$
used in the proof of Lemma \ref{lemma-V-implies-C-general}
which induces by adjunction a map
$$
\epsilon_Y^{-1} Rf_{\tau', *}\epsilon_{X, *}f_{\tau}^{-1}\mathcal{F} \to
Rf_{\tau, *}f_\tau^{-1}\mathcal{F}
$$
Taking trunctions we find a map
$$
\epsilon_Y^{-1}K
\longrightarrow
\tau_{\leq n}Rf_{\tau, *}f_\tau^{-1}\mathcal{F}
$$
which is an isomorphism by Lemma \ref{lemma-V-implies-C-general};
the lemma applies because $f_\tau^{-1}\mathcal{F}$ is in $\mathcal{A}_X$
by Lemma \ref{lemma-A}. Choose a distinguished triangle
$$
\epsilon_{Y, *}\mathcal{F} \to K \to L \to \epsilon_{Y, *}\mathcal{F}[1]
$$
The map $\mathcal{F} \to f_{\tau, *}f_\tau^{-1}\mathcal{F}$
is injective as $\{X \to Y\}$ is a $\tau$-covering. Thus
$\epsilon_{Y, *}\mathcal{F} \to
\epsilon_{Y, *}f_{\tau, *}f_\tau^{-1}\mathcal{F} =
f_{\tau', *}f_{\tau'}^{-1}\epsilon_{Y, *}\mathcal{F}$
is injective too. 
Hence $L$ only has nonzero cohomology sheaves in degrees $0, \ldots, n$.
As $f_{\tau', *}f_{\tau'}^{-1}\epsilon_{Y, *}\mathcal{F}$
is in $\mathcal{A}'_Y$ by (\ref{item-restriction-A}) and
(\ref{item-A-and-P}) we conclude that
$$
H^0(L) = \Coker(\epsilon_{Y, *}\mathcal{F} \to
f_{\tau', *}f_{\tau'}^{-1}\epsilon_{Y, *}\mathcal{F})
$$
is in the weak Serre subcategory $\mathcal{A}'_Y$. For $1 \leq i \leq n$
we see that $H^i(L) = R^if_{\tau', *}f_{\tau'}^{-1}\epsilon_{Y, *}\mathcal{F}$
is in $\mathcal{A}'_Y$ by (\ref{item-restriction-A}) and
(\ref{item-A-and-P}). Pulling back the distinguished triangle
above by $\epsilon_Y$ we get the distinguished triangle
$$
\mathcal{F} \to \tau_{\leq n}Rf_{\tau, *}f_\tau^{-1}\mathcal{F}
\to \epsilon_Y^{-1}L \to \mathcal{F}[1]
$$
Since $\xi$ maps to zero in the middle term we find
that $\xi$ is the image of an element
$\xi' \in H^n_\tau(Y, \epsilon_Y^{-1}L)$.
By Lemma \ref{lemma-V-implies-cohomology-general} we have
$$
H^n_{\tau'}(Y, L) = H^n_\tau(Y, \epsilon_Y^{-1}L),
$$
Thus we may lift $\xi'$ to an element of $H^n_{\tau'}(Y, L)$
and take the boundary into
$H^{n + 1}_{\tau'}(Y, \epsilon_{Y, *}\mathcal{F})$
to see that $\xi$ is in the image of the canonical map
$H^{n + 1}_{\tau'}(Y, \epsilon_{Y, *}\mathcal{F}) \to
H^{n + 1}_\tau(Y, \mathcal{F})$.
By locality of cohomology for
$H^{n + 1}_{\tau'}(Y,\epsilon_{Y, *}\mathcal{F})$, see
Lemma \ref{lemma-kill-cohomology-class-on-covering},
we conclude.
\end{proof}

\end{document}
